El análisis de dispersión entre el ingreso mensual de los clientes y el monto del préstamo no evidencia un patrón o tendencia claramente definida. Los datos se encuentran distribuidos de forma heterogénea a lo largo del gráfico, sin observarse una relación lineal aparente que permita concluir que un mayor nivel de ingresos esté asociado sistemáticamente con préstamos de mayor o menor cuantía. De igual forma, la distribución por sector económico tampoco muestra diferencias visuales que permitan identificar un comportamiento distintivo entre los distintos grupos analizados.
Con base en el análisis exploratorio realizado, la relación entre ingreso, monto del préstamo y sector económico no proporciona evidencia suficiente para considerarse un factor explicativo del comportamiento de la cartera crediticia. En consecuencia, estos resultados no permiten orientar estrategias de gestión del riesgo crediticio sustentadas únicamente en dichas variables, siendo necesario complementar el estudio con análisis estadísticos y variables adicionales que permitan identificar los factores con mayor capacidad explicativa.
Por lo anterior, las políticas de evaluación y administración del riesgo crediticio deberían continuar considerando un enfoque integral que incorpore múltiples variables, tales como el historial de pago, la antigüedad laboral, el nivel de endeudamiento, la capacidad de pago y otros indicadores financieros. La incorporación de estos factores permitirá construir modelos de riesgo más robustos y mejorar la identificación de los principales determinantes de la mora.